\documentclass[12pt,a4paper]{article}
\newcommand{\CadernoDisciplina}{Análise Exploratória de Dados}
\newcommand{\CadernoCorHex}{17365D}
% Preâmbulo compartilhado dos cadernos de exercícios UNIFACOL.
% Definir \CadernoDisciplina e, opcionalmente, \CadernoCorHex antes deste input.
\providecommand{\CadernoDisciplina}{Disciplina}
\providecommand{\CadernoCorHex}{17365D}
\usepackage[a4paper,margin=2.05cm,headheight=15pt]{geometry}
\usepackage[T1]{fontenc}
\usepackage[utf8]{inputenc}
\usepackage[brazil]{babel}
\usepackage{lmodern,microtype,amsmath,amssymb}
\usepackage{enumitem,booktabs,tabularx,array,longtable,adjustbox,multirow}
\usepackage{xcolor,hyperref,graphicx,fancyhdr}
\usepackage[most]{tcolorbox}
\usepackage{tikz}
\usetikzlibrary{arrows.meta,positioning,fit,calc,patterns,shapes.geometric}
\definecolor{disciplinecolor}{HTML}{\CadernoCorHex}
\definecolor{stimulusgrey}{HTML}{EEF2F5}
\definecolor{answerorange}{HTML}{FFF0DE}
\definecolor{answerframe}{HTML}{B85C00}
\definecolor{supportgrey}{HTML}{F7F7F7}
\hypersetup{colorlinks=true,linkcolor=disciplinecolor,urlcolor=disciplinecolor}
\newcolumntype{Y}{>{\raggedright\arraybackslash}X}
\newcolumntype{C}{>{\centering\arraybackslash}X}
\setlength{\parindent}{0pt}
\setlength{\parskip}{.42em}
\setlist[enumerate]{leftmargin=*,itemsep=.28em,topsep=.35em}
\pagestyle{fancy}
\fancyhf{}
\lhead{UNIFACOL --- \CadernoDisciplina}
\rhead{Caderno ENADE --- 2026.2}
\cfoot{\thepage}
\newcommand{\ItemHeader}[4]{%
  \par\medskip\noindent
  \colorbox{disciplinecolor}{\parbox{\dimexpr\textwidth-2\fboxsep\relax}{\color{white}\bfseries Questão #1\hfill #2}}%
  \par\smallskip\noindent\textbf{Competência:} #3\par
  \noindent\textbf{Suporte:} #4\par\smallskip
}
\newcommand{\TextoBase}{\noindent\colorbox{stimulusgrey}{\parbox{\dimexpr\textwidth-2\fboxsep\relax}{\textbf{TEXTO-BASE}}}\par\smallskip}
\newcommand{\Comando}{\par\smallskip\noindent\textbf{COMANDO.} }
\newcommand{\FonteDidatica}{\par\smallskip\noindent{\footnotesize\textit{Fonte: situação e dados elaborados para fins didáticos.}}}
\newcommand{\FonteExterna}[1]{\par\smallskip\noindent{\footnotesize\textit{Fonte: #1}}}
\newcommand{\EspacoDiscursiva}{\par\noindent\rule{\textwidth}{.2pt}\vspace{1.15cm}\par\noindent\rule{\textwidth}{.2pt}}
\newtcolorbox{AnswerBox}{enhanced,breakable,colback=answerorange,colframe=answerframe,boxrule=.7pt,arc=1mm,title=Resposta explicada,fonttitle=\bfseries,coltitle=white,before skip=7pt,after skip=7pt}
\newtcolorbox{DocumentBox}[1][Documento ou artefato]{enhanced,breakable,colback=supportgrey,colframe=black!55,boxrule=.6pt,arc=.7mm,title={#1},fonttitle=\bfseries,coltitle=white,before skip=7pt,after skip=7pt}
\newtcolorbox{MemoBox}{enhanced,breakable,colback=white,colframe=disciplinecolor,boxrule=.7pt,arc=.7mm,title=Memorando,fonttitle=\bfseries,coltitle=white,before skip=7pt,after skip=7pt}
\newcommand{\LegendaDidatica}[1]{\par\smallskip{\footnotesize\textit{#1. Fonte: elaborado para fins didáticos.}}}

\setcounter{secnumdepth}{0}
\begin{document}
\begin{center}
{\Large\bfseries UNIFACOL}\\[2mm]
{\large Avaliação da II Unidade}\\[1mm]
{\Large\bfseries VERSÃO B}
\end{center}
\begin{DocumentBox}[Identificação]
\begin{tabularx}{\textwidth}{@{}p{.48\textwidth}p{.48\textwidth}@{}}
Estudante: \hrulefill & Turma: \hrulefill\\[3mm]
Professor: Cayo Oliveira & Data: \rule{2.7cm}{.2pt}/2026
\end{tabularx}
\end{DocumentBox}
\textbf{Instruções.} Esta avaliação contém seis questões objetivas e duas discursivas. Nas objetivas, assinale uma única alternativa. Nas discursivas, mostre cálculos intermediários e justifique decisões com as evidências do texto-base. Respostas sem desenvolvimento podem não receber pontuação integral.
\bigskip
% origem: AED-C14-O05
\ItemHeader{1}{Objetiva}{Distinguir alucinação de desatualização}{Política revogada}
\TextoBase
O recuperador trouxe uma política revogada, marcada como versão vigente por erro de metadado. O modelo citou fielmente o trecho e não inventou conteúdo, mas a decisão contrariou a norma atual.
A resposta usou corretamente um documento que deixou de valer; a avaliação deve separar recuperação desatualizada de conteúdo sem sustentação.
\FonteDidatica
\Comando A classificação mais precisa é
\begin{enumerate}[label=\Alph*)]
  \item falha de custo, porque versões antigas usam mais tokens.
  \item falha de vigência/proveniência na recuperação; fidelidade ao contexto não garantiu correção atual.
  \item sucesso, porque a citação existe.
  \item alucinação textual exclusiva, porque toda resposta incorreta é alucinação.
  \item falha exclusiva de estilo, resolvida com prompt mais formal.
\end{enumerate}

% origem: AED-C13-O01
\ItemHeader{2}{Objetiva}{Distinguir responsabilidades agentic}{Relatório mensal recorrente}
\TextoBase
Um fluxo recebe a solicitação, consulta métricas, compara o mês anterior, produz narrativa e pede aprovação humana. A equipe chama de ``agente'' tanto o modelo quanto uma função SQL e um manual de procedimento.
O relatório alimenta decisões de estoque, e cada etapa deve possuir entrada, saída, permissão e condição de parada verificáveis.
\FonteDidatica
\Comando A descrição tecnicamente adequada é
\begin{enumerate}[label=\Alph*)]
  \item o agente coordena estado e decisões; SQL é ferramenta; o manual pode constituir skill ou playbook; o LLM apoia interpretação.
  \item o LLM é a ferramenta, o SQL é o agente e o manual é a memória episódica.
  \item todo componente é um agente porque participa do mesmo fluxo.
  \item a skill executa diretamente qualquer ação e torna permissões desnecessárias.
  \item a memória decide o próximo passo e o agente apenas armazena texto.
\end{enumerate}

% origem: AED-C10-D13
\ItemHeader{3}{Discursiva}{Integrar evidências e justificar decisão}{Caso profissional do capítulo}

\TextoBase

As áreas comercial e financeira usam o nome ``receita líquida'', porém uma exclui apenas cancelamentos e a outra também exclui impostos. O agente responde de formas diferentes conforme as colunas disponíveis, sem informar versão, responsável ou grão da definição.

\FonteDidatica

\Comando Proponha uma entrada de catálogo para ``receita líquida'' contendo fórmula, grão, dimensões permitidas, filtros, responsável, vigência, versão e três testes. Explique como o agente deve tratar a ambiguidade entre as duas definições.
\EspacoDiscursiva

% origem: AED-C09-O03
\ItemHeader{4}{Objetiva}{validação}{Situação profissional e suporte quantitativo}

\TextoBase

O assistant afirma queda de 18\%, mas a tabela de referência contém 120 antes e 110 depois. A resposta será enviada à diretoria. Antes, a aplicação precisa recalcular a variação relativa e tratar a frase do modelo como hipótese sujeita à fonte.

\FonteDidatica

\Comando Recalcule a variação observada e selecione a ação de validação antes da publicação.

\begin{enumerate}[label=\Alph*)]
  \item aumentar contexto sem conferir
  \item trocar system por user
  \item aceitar 18\% por ser resposta do modelo
  \item reduzir temperatura e publicar
  \item recalcular: queda de 8,33\%; bloquear ou corrigir resposta e registrar divergência
\end{enumerate}

% origem: AED-C12-O07
\ItemHeader{5}{Objetiva}{Garantir saída estruturada}{Gráfico sugerido pelo Llama}
\TextoBase
O modelo deve sugerir tipo de gráfico, campos e justificativa. Em testes, produz nomes de colunas inexistentes e ocasionalmente inclui código. A interface espera um objeto estruturado.
A especificação do gráfico será consumida por componente separado, exigindo campos permitidos, tipos verificáveis e rejeição de propriedades inesperadas.
\FonteDidatica
\Comando O fluxo de validação deve
\begin{enumerate}[label=\Alph*)]
  \item inserir todas as colunas do DataFrame no prompt até o modelo parar de errar.
  \item validar esquema, lista permitida de gráficos, existência e tipos das colunas; rejeitar ou reparar de modo controlado antes de renderizar.
  \item remover justificativa e validação para diminuir tokens.
  \item extrair nomes por expressão regular e executar eventual código junto da visualização.
  \item aceitar qualquer JSON sintaticamente válido e renderizar os campos recebidos.
\end{enumerate}

% origem: AED-C11-O03
\ItemHeader{6}{Objetiva}{Distinguir modelo semântico de visualização}{Dois totais para a mesma receita}
\TextoBase
No piloto, Looker com Gemini consulta uma medida certificada que exclui cancelamentos e usa câmbio do fechamento. Um app Streamlit soma a coluna bruta e usa câmbio do dia. Ambos desenham gráficos corretos para seus respectivos resultados, mas os totais divergem.
A divergência afetará o fechamento financeiro, por isso uma mudança apenas visual não resolve o desacordo entre contratos de cálculo.
\FonteDidatica
\Comando O diagnóstico e a primeira correção adequados são
\begin{enumerate}[label=\Alph*)]
  \item alinhar definição, grão, filtros e regra cambial em uma camada semântica versionada antes de comparar interfaces.
  \item aumentar o modelo de linguagem para que escolha o total mais plausível.
  \item calcular a média dos dois totais e publicá-la como conciliação.
  \item manter ambos os números, pois ferramentas diferentes necessariamente produzem verdades diferentes.
  \item atribuir a divergência ao tipo de gráfico e padronizar as cores.
\end{enumerate}

% origem: AED-C10-O01
\ItemHeader{7}{Objetiva}{semântica}{Situação profissional e suporte quantitativo}

\TextoBase

A métrica ``receita líquida'' possui contrato oficial que exclui cancelamentos e impostos, com vigência e responsável. O usuário não conhece a fórmula e pede apenas ``receita''. O agente precisa resolver a intenção pela camada semântica, não inventar cálculo a partir de nomes de colunas. O catálogo registra ainda moeda, data de competência e dimensão de canal permitida. Como o relatório será comparado ao fechamento financeiro, uma soma sintaticamente válida sobre coluna homônima seria rejeitada mesmo que produzisse um número plausível. O pedido abrange apenas lojas ativas no Nordeste e o segundo trimestre; esses filtros também pertencem ao plano semântico e devem aparecer na explicação. Antes da execução, o agente precisa mostrar qual métrica governada, quais dimensões e qual janela temporal resolveu, permitindo que o analista confirme a interpretação.

Essa confirmação deve ocorrer antes do fechamento financeiro.

\FonteDidatica

\Comando Para responder sem criar uma definição paralela de receita, o agente deve

\begin{enumerate}[label=\Alph*)]
  \item duplicar definição em cada prompt
  \item usar a maior medida
  \item pedir ao LLM que invente a fórmula
  \item resolver a expressão pela métrica governada, não pela soma inferida de uma coluna homônima
  \item somar toda coluna receita
\end{enumerate}

% origem: AED-C11-D13
\ItemHeader{8}{Discursiva}{Desenhar arquitetura híbrida}{BI corporativo e Streamlit}
\TextoBase
O time quer preservar a camada semântica corporativa, mas precisa de simulação customizada, comentários humanos e documento final assinado. Hoje o app usa credencial administrativa e copia dados para CSV.
O documento assinado gera efeito administrativo, de modo que cada transição de estado deve conservar autor, horário e evidência consultada.
\FonteDidatica
\Comando Desenhe uma arquitetura corrigida, incluindo identidade do usuário, consulta governada, política por linha, serviço de simulação, estado de revisão, evidências, assinatura, logs minimizados e critérios de teste e rollback.
\EspacoDiscursiva

\end{document}
